\section{Conclusion}
\label{sec:conclusion}

We presented ALM-MIL for examination-level multimodal multi-label recognition
from ordered image sequences and paired findings reports. ACPE uses acquisition
order as a structural anchor and adapts contextual coordinates to visual
transitions, while injecting positional information through complementary
absolute and relative routes. LCCF organizes multimodal evidence around
individual labels through label-wise visual selection, coexistence-aware
hypergraph reasoning, and label-conditioned textual evidence integration with
adaptive fusion. Experiments on four clinical cohorts and four public datasets
demonstrate the effectiveness of the framework and its two core components.
Overall, the results show that preserving acquisition structure and organizing
multimodal evidence at the label level are complementary requirements for
reliable examination-level multi-label recognition.

\cntranslation{%
本文提出 ALM-MIL，用于基于有序图像序列和配对所见报告的检查级多模态多标签识别。ACPE 以采集顺序为结构参照，并根据视觉变化自适应构建上下文坐标，通过互补的绝对路径和相对路径注入位置信息。LCCF 则围绕各个标签组织多模态证据，通过标签级视觉选择、共存感知超图推理以及带自适应融合的标签条件文本证据整合，实现标签级的图像与文本建模。在四个临床队列和四个公开数据集上的实验验证了该框架及其两个核心模块的有效性。总体而言，实验结果表明，保留采集结构并在标签层面组织多模态证据，是实现可靠检查级多标签识别的两个互补条件。
}
